\section{Induction}
The rest of the deductive rules of Type Theory go on to give schemas to construct, manipulate, and
eliminate various constructions (known as type formers/inductive types). A number of these
constructions are \textit{synthetically} baked into the deductive system, known as \textit{type formers}.
Type Theories also allow the user to define their own constructions, known as inductive types,
requiring the same criteria to be defined. 

The general scheme for these type formers/inductive types is as follows. An type former/inductive type is defined by:
\begin{itemize}
  \item A list of \textbf{\textit{formation rules}}, describing how the inductive type can be
    introduced into context.
\item A list of \textbf{\textit{constructors/indroduction rules}}, to introduce new terms of these types.
\item \textbf{\textit{Elimination rules}} to give a recipe how to use/consume elements of the type. \\
    (Usually, this involves a \textit{recursion principle} to define functions on the type, and a
    stronger \textit{induction principle} to define dependent functions on the type.)
\item A \textbf{\textit{computation rule}}, that tells how the elimination acts on the constructors. 
\item An (optional) \textbf{\textit{uniqueness principle}}, dictating how each element of the type
  is uniquely determined by the result of elimination rules applied to it.
\end{itemize}

Here is a brief, informal description of these type formmers in Type Theory. Each of them also have a logical interpretation
to them under `propositions as types'. A detailed, formal description
is avaiable in Appendix I of \cite{HoTT}.

\subsection{Some Type Formers/Inductives}
\subsubsection{Function Types}
Given $A, B : \U$, we can form $A \to B : \U$, the type of functions from $A$ to $B$,
syntactically given by a simple $\lambda$-calculus. Introducing a function can be done by specifying a 
formula as $(\lambda x.\Phi)$ where $\Phi$ is a formula in the variable $x$. Elimination is performed by
evaluation, computed as $(\lambda x.\Phi)(a) \deq \Phi[a/x]$, replacing all occurences of $x$ with $a$.
We get $f(a) : B$ for any $f : A \to B, a : A$.

Iterated function types such as $A \to (B \to C))$ are effectively multivarible function types (this is called currying, due to Product-Exponential adjunction),
so the parentheses are omitted and associated by right from default, like $A \to B \to C$.
Currying is also used in lambda notation, writing $(\lambda x.\lambda y.\Phi)$ instead of $(\lambda x.(\lambda y.\Phi))$.

Logically, $A \to B$ correponds to $A \implies B$.

\subsubsection{Dependent Function Types ($\Pi$-types)}
Given $A : \U, B : A \to \U$, the type of dependent functions is $\Prod{x:A} B(x)$, functions with
codomain dependent on the domain. They follow exactly the same rules as regular function types, but $f(a) : B(a)$
for $a : A, f : \Prod{a:A}{B(a)}$.

Logically, $\Pi$-types correspond to for-all predicates ($\forall$).

\subsubsection{Product Types}
Given $A, B : \U$, the product type $A \times B : \U$ consists of pairs $(a, b)$, given by $a : A, b : B$. 
Let $E : A \times B \to \U$, then to define $f : \Prod{z:A\times B}{E(z)}$, it suffices to define $f$
for the special case where $x \deq (a,b)$ for some $a : A, b : B$. This is the \textit{induction principle}, denoted
$\mathsf{ind}_{A\times B}$. The computation rule enforces the defined function matches in the special case.
Formally, this is
\[
  \mathsf{ind}_{A \times B} : \Prod{E : A \times B \to \U}{\left(\Prod{x:A}{\Prod{y:B} E((x,y))} \to \Prod{z:A \times B} E(z)\right)}
\]
with the computation rule
\[
  \mathsf{ind}_{A \times B} (E, g, (a,b)) \deq g(a)(b)
\]
But from now on we'll often condense this into informal language for convenience.

Logically, $A \times B$ corresponds to $A \land B$.

\subsubsection{Coproduct Types}
Given $A, B : \U$, the coproduct type $A + B$ is given by two generators $\inl : A \to A + B$, and $\inr : B \to A + B$.
For its induction principle, to define a function $f : \Prod{z : A + B} E(z)$, it suffices to define $f$ for the 
special case where $z \deq \inl (x)$ and the case where $z \deq \inr (y)$.

Logically, $A + B$ corresponds to $ A \lor B$.

\subsubsection{Dependent Pair Types ($\Sigma$-types)}
This is the dependent version of a product type. Given $A : \U, B : A \to \U$, we have $\Sum{a:A}{B(a)} : \U$.
Say $a : A, b : B(a)$, one can form $(a,b) : \Sum{a:A}{B(a)}$. 
Let $E : \Sum{a:A}{B(a)} \to \U$. To define $f : \Prod{c:\Sum{a:A}{B(a)}} {E(c)}$, it suffices to define
$f$ in the special case where $c \deq (a,b)$ for some $a : A, b : B(a)$.

Logically, $\Sigma$-types correspond to there-exists predicate $(\exists)$.

\subsubsection{0, 1 and 2}
The empty type $\0$ is the type with no constructors, and the induction principle, that to define
$f : \Prod{x:\0}{E(x)}$, it suffices to \textit{not define anything}. By default, we have a function,
$\Prod{x:\0}{E(x)}$.

The unit type $\1$ is the type with a single constructor, $\star : \1$. The induction principle says that
to define $f : \Prod{x:\1}{E(x)}$, it suffices to define $f$ for $x \deq \star$.

The boolean type $\2$ is the type with two constructors, $0_2 : \2$ and $1_2 : \2$. To define 
$f : \Prod{x:\2}{E(x)}$, it suffices to define for $x \deq 0_2$ and $x \deq 1_2$.
Alternatively, one can define $\2 :\deq \1 + \1$.

Logically, $\0$ is a contradiction ($\bot$) and $\1$ is True ($\top$).

\subsubsection{Natural Numbers Type ($\N$)}
The natural numbers type is the prototypical example of an \textit{Inductive Type}. It is generated by the constructors
\begin{itemize}
  \item $\mathsf{zero} : \N$
  \item $\mathsf{succ} : \N \to \N$
\end{itemize}
Let $E : \N \to \U$. Then to define $E(n)$ for all $n: \N$, it suffices to define it for the following special case 
$n \deq \mathsf{zero}$, and for all $m : \N$, given $f(m)$ define it for $n \deq \mathsf{succ}(m)$.
(This is almost the usual principle of mathematical induction!)

\subsubsection{Identity Types}
Given $A : \U$, we have the identity type family $\Id_A : A \to A \to \U$. Given $x, y: A$, we have $\Id_A(x,y)$, also denoted as
$x =_A y$, which is the type of all \textit{propositional equalities} between $x$ and $y$. The introduction rule is by using
$\refl : \Prod{a:A} (a =_A a)$. Identities being types themselves allows equality of objects to also 
be treated as points, susceptible to comparison with other equalities, and thus proofs of equalities as mathematical objects susceptible to proof of equality, ad infinitum.
Thus, one obtains an infinite-order logic as opposed to set theory's first-order.

The induction principle for Identity Types is called path induction, which says 
given $C : \Prod{x, y : A} (x =_A y) \to \U$, a function $c : \Prod{x : A}{C(x,x,\refl_x)},$
there is a function $f : \Prod{x,y : A}{\Prod{p: x =_A y} C(x,y,p)}$ such that $f(x,x,\refl_x) :\deq c(x)$.
That is, to define a function on paths between points in $A$, it suffices to exhibit a function on all
reflexive loops $\refl_x$ over all points $x : A$.

\subsection{Custom Inductive Types}
Type Theories provide a schema to build newer, user-defined inductive types as well.
Conceptually, an inductive type is a type \textit{freely generated} by a number of constructors, which may involve functions
that output the given type by (possibly) taking as input terms of the same type. This intuition needs
to be formalized carefully so as to ensure that only meaningful inductive types may be constructed.

\textit{\textbf{Remark.} Inductive types are objects one can perform induction on! 
Similarly to $\N$, the \textit{induction principle} for a given inductive type is the statement that to prove a
statement for all its instances, it suffices to prove it in a few special cases, built on the constructors
of the type.}

\subsubsection{W-Types}
There is a complete formalism of inductive types up to certain demands for the constructors.

\begin{definition}[W-Types]
  Given $A : \U, B : A \to \U$, the W-Type $\W{a}{A}{B}$ is defined by the constructor
  \[
    \mathsf{sup} : \left(\Prod{a:A}{B(a) \to \W{a}{A}{B}} \right) \to \W{x}{A}{B}
  \]
  and the induction principle, to prove $E : \w \to \U$ for all elements of $\w$, it suffices
  to prove it for $\mathsf{sup}(a,f)$ for each $a : A$, and $f : (B(a) \to \w)$, assuming it holds
  for all $f(b)$ with $b : B(a)$.
\end{definition}

A W-Type $\w$ may be interpreted as the type of well-founded trees with $A$-many ways to branch out from
each node, corresponding to a constructor.
Each $a : A$ corresponds to a constructor, each of which take $B(a)$-many instances of $\w$
as branches. The type $\w$ is the type of objects \textit{freely generated} by these constructors.
If an inductive type may be expressed equivalently as generated by $\mathsf{sup}$ as described, it can be formalized as a W-Type.

\begin{itemize}
  \item The boolean type $\2$ is equivalent to the W-Type $\w$ where 
    $A \deq \2$ with $B(0_2) \deq B(1_2) \deq \0$, that is generated by 
    two constructors, $0_2 : \2$ and $1_2 : \2$.
  \item The natural numbers type $\N$ is equivalent to $\w$ where $A \deq \2$,
    with $B(0_2) \deq \0$ and $B(1_2) \deq \1$, i.e., generated by two constructors,
    corresponding to $\mathsf{zero : \N}$, and $\mathsf{succ} : \N \to \N$.
  \item The type $\W{x}{\1+A}{B}$ with $B(\inl (\star)) \deq \0$ and $B(\inr (a)) \deq \1$,
    corresponds to the type $\mathsf{List}(A)$ of finite length strings over the alphabet $A$, i.e.
    generated by $\mathsf{nil} : \mathsf {List}(A)$ and $\mathsf{cons} : A \to \mathsf{List}(A) \to \mathsf{List}(A)$.
\end{itemize}
